\section{Síntesis y ampliación}

\subsection{Resumen de la clase}

\begin{figure}[H]
\centering
\includegraphics[width=0.95\textwidth]{slides-images/slide-58.jpg}
\caption{Resumen de Deep RL: tres bloques -- conceptos básicos, Q-Learning y Policy Gradient\protect\footnotemark}
\end{figure}
\footnotetext{Diapositivas, página 58.}

Esta clase presentó de manera sistemática el contenido central del Deep Reinforcement Learning, dividido principalmente en tres partes:

\textbf{Uno. Conceptos básicos de RL}
\begin{itemize}
    \item El Agent interactúa con el Environment mediante Action
    \item Los pares State-Action constituyen los datos básicos del aprendizaje
    \item El objetivo es maximizar el retorno acumulado con descuento $R_t = \sum \gamma^i r_i$
    \item La Q-function $Q(s,a)$ mide el retorno futuro esperado de un par estado-acción
\end{itemize}

\textbf{Dos. Value Learning / Q-Learning}
\begin{itemize}
    \item Se usa DQN para aproximar la Q-function
    \item La política se deriva de forma determinista mediante $\arg\max_a Q(s,a)$
    \item Aplicable a espacios de acción discretos
    \item Logró un desempeño superior al humano en juegos de Atari
\end{itemize}

\textbf{Tres. Policy Gradient}
\begin{itemize}
    \item Aprende y optimiza directamente la función de política $\pi(s)$
    \item Aplicable a espacios de acción continuos
    \item Se entrena mediante $\text{loss} = -\log P(a_t|s_t) \cdot R_t$
    \item Logró avances en aplicaciones complejas como la conducción autónoma y el Go
\end{itemize}

\subsection{Lecturas adicionales}

\begin{itemize}
    \item \textbf{Artículo original de DQN}: Mnih et al., ``Human-level control through deep reinforcement learning,'' \textit{Nature}, 2015
    \item \textbf{Algoritmo REINFORCE}: Williams, ``Simple statistical gradient-following algorithms for connectionist reinforcement learning,'' \textit{Machine Learning}, 1992
    \item \textbf{AlphaGo}: Silver et al., ``Mastering the game of Go with deep neural networks and tree search,'' \textit{Nature}, 2016
    \item \textbf{AlphaZero}: Silver et al., ``A general reinforcement learning algorithm that masters chess, shogi, and Go through self-play,'' \textit{Science}, 2018
    \item \textbf{Simulador VISTA}: Amini et al., ``Learning Robust Control Policies for End-to-End Autonomous Driving From Data-Driven Simulation,'' \textit{IEEE RA-L}, 2020
    \item \textbf{RLHF}: Ouyang et al., ``Training language models to follow instructions with human feedback,'' \textit{NeurIPS}, 2022
    \item \textbf{PPO}: Schulman et al., ``Proximal Policy Optimization Algorithms,'' 2017------la variante de Policy Gradient más popular actualmente
    \item \textbf{Métodos Actor-Critic}: combinan las ventajas de Value Learning y Policy Gradient, aprendiendo simultáneamente la función de valor y la función de política
    \item \textbf{Sitio del curso}: \url{https://introtodeeplearning.com}
\end{itemize}

\begin{knowledgebox}{El futuro del Deep RL}
El aprendizaje por refuerzo profundo sigue siendo un campo en rápido desarrollo. Los focos de investigación actuales incluyen:
\begin{itemize}
    \item \textbf{Sample Efficiency}: cómo alcanzar el mismo desempeño con menos datos de interacción
    \item \textbf{Multi-Agent RL}: escenarios en los que varios Agent aprenden y colaboran o compiten simultáneamente
    \item \textbf{Offline RL}: entrenamiento sin conexión a partir de datos históricos existentes, evitando el costo y el riesgo de la interacción en línea
    \item \textbf{Foundation Models + RL}: incorporar el conocimiento de los modelos preentrenados a gran escala al RL para acelerar el aprendizaje
    \item \textbf{RLHF y AI Alignment}: entrenar sistemas de IA más seguros y más útiles mediante retroalimentación humana
\end{itemize}
\end{knowledgebox}


%% --- Fin del contenido --- %%

\end{document}
